% !TEX program = xelatex
\documentclass[12pt, a4paper]{ctexart}
\usepackage{amsmath}
\usepackage{amssymb}
\usepackage{geometry}
\geometry{left=2.5cm,right=2.5cm,top=2.5cm,bottom=2.5cm}
\usepackage{listings}
\usepackage{xcolor}
\usepackage{draftwatermark}

% 水印配置
\SetWatermarkText{贺昱超学习资料}
\SetWatermarkLightness{0.92}
\SetWatermarkScale{1.0}
\SetWatermarkAngle{45}

% 代码高亮配置
\lstset{
    language=C++,
    basicstyle=\ttfamily\small,
    keywordstyle=\color{blue}\bfseries,
    commentstyle=\color{green!50!black},
    stringstyle=\color{purple},
    numbers=left,
    numberstyle=\tiny\color{gray},
    breaklines=true,
    frame=single,
    backgroundcolor=\color{gray!5},
    showstringspaces=false
}

\title{\textbf{特征点提取与光流法追踪}}

\date{2026年8月}

\begin{document}

\maketitle

\section{引言：数学的美妙巧合}
在计算机视觉中，我们要完成两件事：一是找出画面里最容易追踪的“特征点”（角点提取），二是计算这些点在两帧之间的移动速度（光流法）。
在调用封装接口时，这看似是两个毫不相干的函数。但当我们今天用 C++ 纯手工实现它们底层时，你会发现一个令人惊叹的数学巧合：\textbf{这两个算法的核心，竟然是求解完全相同的同一个矩阵！}

\section{核心基石：结构张量矩阵 (Structure Tensor)}
无论是角点检测，还是光流法，第一步都是计算图像的微积分导数：水平偏导 $I_x$ 和垂直偏导 $I_y$。
在一个像素点周围的 $W \times W$ 邻域内，我们构建一个 $2 \times 2$ 的矩阵 $M$，称为\textbf{结构张量}：
\begin{equation}
M = 
\begin{bmatrix}
\sum I_x^2 & \sum I_x I_y \\
\sum I_x I_y & \sum I_y^2
\end{bmatrix}
\end{equation}
这便是连接感知与运动的终极钥匙。

\section{特征点提取 (Shi-Tomasi)：求矩阵的特征值}
\subsection{原理解析}
什么样的点才是“好点”（角点）？如果一个点在图像中是平坦的墙面，矩阵 $M$ 的值全接近 $0$；如果是一条垂直的边缘，它在 Y 方向没有变化。
在线性代数中，矩阵的\textbf{特征值 (Eigenvalues)} 代表了数据分布在主成分上的方差。Shi 和 Tomasi 提出：\textbf{只要矩阵 $M$ 的两个特征值 $\lambda_1, \lambda_2$ 中的极小值足够大，说明这个点在任何方向都有剧烈梯度，它必定是一个角点！}

\subsection{C++ 硬核求解}
对于 $2 \times 2$ 矩阵 $M = \begin{bmatrix} a & c \\ c & b \end{bmatrix}$，求其特征值即解一元二次方程 $\lambda^2 - (a+b)\lambda + (ab-c^2) = 0$。我们在代码中直接手写求根公式得出最小特征值：
\begin{lstlisting}
float a = Sxx.at<float>(y, x); // ΣIxx
float b = Syy.at<float>(y, x); // ΣIyy
float c = Sxy.at<float>(y, x); // ΣIxy

// 计算最小特征值 λ_min
float lambda_min = (a + b - sqrt((a - b)*(a - b) + 4 * c * c)) / 2.0f;
\end{lstlisting}
通过设定阈值并进行非极大值抑制（找出局部最亮的点），我们就用纯数学抠出了视野中所有的角点。

\section{光流法 (Lucas-Kanade)：解矩阵的方程组}
\subsection{原理解析}
找到了锚点后，我们要计算它的速度向量 $\mathbf{v} = [u, v]^T$。根据光照不变性定理与一阶泰勒展开，我们能推导出超定方程组 $A\mathbf{v} = -b$。
通过最小二乘法，两边同乘 $A^T$，左边得到的矩阵 $A^T A$，\textbf{正是我们刚才寻找角点时构建的结构张量矩阵 $M$！}
方程组最终化简为：
\begin{equation}
\begin{bmatrix}
\sum I_x^2 & \sum I_x I_y \\
\sum I_x I_y & \sum I_y^2
\end{bmatrix}
\begin{bmatrix} u \\ v \end{bmatrix}
=
\begin{bmatrix}
-\sum I_x I_t \\
-\sum I_y I_t
\end{bmatrix}
\end{equation}

\subsection{C++ 克莱姆法则求解}
既然矩阵相同，光流法实际上就是在解以该矩阵为系数的二元一次方程组。我们在代码中计算矩阵的行列式 $det$，并用克莱姆法则直接求解出屏幕上那一抹赛博朋克的尾迹：
\begin{lstlisting}
// 1. 计算行列式 det = ad - bc
double det = sum_Ixx * sum_Iyy - sum_Ixy * sum_Ixy;

// 物理意义：行列式为0代表特征值极小（白墙），方程无解，抛弃该点
if (abs(det) < 1e-5) return false; 

// 2. 使用克莱姆法则求解瞬时速度 u, v
double u = (sum_Iyy * (-sum_Ixt) - sum_Ixy * (-sum_Iyt)) / det;
double v = (sum_Ixx * (-sum_Iyt) - sum_Ixy * (-sum_Ixt)) / det;
\end{lstlisting}


\end{document}